%---------------------------------------------------------------------
\chapter{Experimentation}
%---------------------------------------------------------------------
The project ran a chain of experiments, each motivated by the previous result and recorded in an
append-only experiment log (\texttt{docs/checkpoints/mg2/EXPERIMENT\_LOG.md}). This chapter
summarises the chain and then reports the main controlled experiment, MG2, and the held-out run.

\section{The Experimental Progression}
Table~\ref{tab:progression} lists the steps. The early steps (MG1) changed one thing at a time on the
single agent; their main finding was that the gains came from the deterministic engine, not from the
model, which motivated the multi-agent hypothesis.

\begin{table}[htbp]
\centering
\caption{Hypothesis, change, result and decision at each step}
\label{tab:progression}
\small
\begin{tabularx}{\linewidth}{@{}l >{\raggedright\arraybackslash}X >{\raggedright\arraybackslash}X >{\raggedright\arraybackslash}X@{}}
\toprule
\textbf{Log} & \textbf{Change} & \textbf{Result} & \textbf{Decision} \\ \midrule
1 & Prompts v8--v12, each an ablation of v7 & v7 best at 69 of 128 trials; the model contributed 3 findings & Stop prompt tuning; find why cases fail \\
2 & Trace every failing case & Seven cases fail every run for deterministic causes & Fix the engine and gate \\
3--4 & Family checks, contradiction gate, one case fix & 30, then 32 of 32; model proposed one finding, rejected & Check on checking-8 \\
5 & Model-free controls on checking-8 & V0, SC, V1 and V2 all 7 of 8; model share 0 & Give the model a measurable role \\
6 & Page loss reported as F15 & 32 of 32 on both sets & Keep \\
7 & All family checks off; thinking prompt & 0 of 21 for v7 and SC; thinking 1 of 21 with false claims & v7 stays; the model never read text \\
9--11 & Retriever and Classifier designed, pre-registered, built & Scripted tests only & Run on the real model \\
12--14 & MG2 on T4: 18 sessions, 1{,}008 trials & No multi-agent arm passes the gate & Nothing selected \\
15 & Stop prompt iteration; declare held-out run & -- & Run held-out once \\
\bottomrule
\end{tabularx}
\end{table}

\section{Hypothesis}
\emph{A Retriever that reads and grades document text, together with a Classifier that maps verified
claims and computed facts to fault families, recovers judgement faults in the ablated arena that the
single-agent v7 loop and the scripted control miss, without adding false findings when the
deterministic checks are on.} The hypothesis is falsified if no multi-agent arm both beats SC in the
ablated arena and passes the rules-on gate (32 of 32 with no false positive).

\section{Experimental Setup}
The arms, arenas, metrics and selection rule were fixed in a pre-registration before the first run:
\begin{itemize}[itemsep=1pt]
\item \textbf{Arms.} D0/SC (controls); S1, the single-agent v7 loop (V1 and V2); M1, the Retriever
followed by a deterministic claim-to-family lookup; M2, the Retriever and the Classifier; M2-nr,
the Classifier alone. M1 and M2-nr isolate the contribution of each role.
\item \textbf{Prompts.} At most four versions declared in advance, each one change from
\texttt{mg2\_agents\_v1}: v2 drops the Retriever's examples, v3 drops the Classifier's family
meanings, v4 left open for one change a result motivates.
\item \textbf{Held fixed.} Model (Qwen3.5-9B Q8\_0), server build, three parallel slots, greedy
decoding, budgets, corpus, cases (tuning-16), two repeats per case, and the scorer.
\item \textbf{Selection rule.} Fewest false positives, then ablated recall, then calls; the chosen
arm must pass the rules-on gate. Only then is a freeze manifest created and the held-out set run.
\item \textbf{Where.} The PC became unavailable, so the 18 development sessions ran on free Kaggle T4
GPUs as a separate series with one process per trial instead of one container (a recorded
deviation). The answers stayed on the laptop, where every session was scored.
\end{itemize}

\section{Results}
Figure~\ref{fig:mg2plot} plots the two quantities that decide the hypothesis for every arm: true positives
recovered in the ablated arena (the benefit) and accepted false positives with the rules on (the
cost). Tables~\ref{tab:mg2-ablated} and~\ref{tab:mg2-rules-on} give the full results.

\begin{figure}[htbp]
\centering
\begin{tikzpicture}
\pgfplotstableread[col sep=comma]{generated/mg2-plot.csv}\mgtwo
\begin{axis}[ybar, bar width=7pt, width=\linewidth, height=6.2cm,
  ymin=0, ylabel={Findings (32 trials per arm)},
  xtick=data, xticklabels from table={\mgtwo}{label},
  x tick label style={rotate=35,anchor=east,font=\scriptsize},
  legend style={at={(0.5,1.03)},anchor=south,font=\scriptsize,draw=none},
  legend columns=2,
  nodes near coords, nodes near coords style={font=\tiny},
  axis lines*=left, ymajorgrids, grid style={gray!25}, enlarge x limits=0.07]
\addplot[fill=BrickRed!70,draw=BrickRed] table[x expr=\coordindex, y=rules_on_fp]{\mgtwo};
\addplot[fill=ForestGreen!60,draw=ForestGreen] table[x expr=\coordindex, y=ablated_tp]{\mgtwo};
\legend{{Accepted false positives (rules on)},{True positives (ablated)}}
\end{axis}
\end{tikzpicture}

\caption{Cost and benefit of each arm in the MG2 development series}
\label{fig:mg2plot}
\end{figure}

No arm supports the hypothesis. The best ablated recall, 6 true positives (out of 42 expected
findings over 32 trials), came from M2 v2 and M1 v1, which with the rules on accepted 25 and 37 false
positives. The Classifier alone made no false positive in the ablated arena and only one or two with
the rules on, but found 2 or 3 findings when the rules were off. Changing the prompt (v2, v3) moved
the counts without changing the pattern. The single-agent arms were the only ones without false
positives, and they found nothing when the rules were off.

\section{Discussion and Decision}
The hypothesis was not supported. The Retriever does read text and does produce claims, which the
single-agent loop never did, but at this model size most of its claims attach a plausible quote to
the wrong family, and the gate, which checks that evidence exists and is not contradicted, lets them
through. The Classifier alone is safe mainly because the gate rejects almost everything it proposes.

By the selection rule no arm could be selected, and on 11 November 2026 the team decided to stop
prompt iteration rather than write a fourth version (log entry 15): three versions across three
multi-agent arms had produced the same false-positive families, and the remaining budget was better
spent on one honest held-out run. The held-out run was declared before it started, with two recorded
deviations: it uses the single-agent arms (the primary V1 versus V2 comparison of the evaluation
protocol) with the controls, plus the Classifier alone with prompt v3 as the arm that ranks first by
the selection rule, reported descriptively; and it runs on a Kaggle T4 with the 24 visible dossiers
uploaded once to a private dataset, verified against two freeze manifests that pin the commit,
protocol, prompts and dossier bytes.

\section{Held-out A/B Result}
The held-out run is reported in Table~\ref{tab:heldout} (Chapter~7). The primary question is whether
V2, the adaptive loop, differs from V1, the fixed sequence, on cases neither was tuned on; the
controls show what the deterministic engine achieves alone on the same cases.
